%%%% CAPÍTULO 6 - RESULTADOS E DISCUSSÃO

\chapter{Resultados e discussão}\label{cap:resultados}

Este capítulo apresenta três conjuntos de evidências: execução funcional na placa, testes automatizados e uma campanha quantitativa. Os conjuntos respondem a perguntas diferentes. A execução em placa mostra que os componentes foram integrados; os testes exercitam casos especificados; e a campanha caracteriza tempo e vazão na configuração registrada. Nenhum deles avalia lucratividade, envio de ordens ou operação em bolsa.

\section{Execução funcional na DE10-Nano}\label{subsec:execucaoFuncionalDe10}

No ensaio preservado pelos autores, o receptor foi iniciado com base MMIO \codigo{0xff200000} e reconheceu assinatura \codigo{0x48465431}, versão 1, profundidade 64, oito palavras por posição e base RX \codigo{0x900}. Depois da conexão ao gerador local, eventos produziram respostas com a mesma sequência.

O \autoref{codigo:saidaFpgaArm} contém três observações desse ensaio. Nos dois primeiros casos, a diferença entre as quantidades do topo excede 500 e a regra do \autoref{codigo:pseudocodigoEstrategia} retorna BUY. No terceiro, a diferença é 121 e a saída é NOOP. O trecho evidencia correlação e coerência com a regra; por ser uma amostra curta, não mede taxa de erro nem cobre todos os eventos possíveis.

\begin{sourcecode}[htb]
\caption{Trecho da saída do receptor com respostas FPGA--ARM}
\label{codigo:saidaFpgaArm}
\begin{lstlisting}[frame=single]
FPGA MMIO bridge enabled: base=0xff200000 span=0x2000
magic=0x48465431 version=1 tx_depth=64 rx_depth=64
slot_words=8 rx_base=0x900 mode=new
Connected to 127.0.0.1:9001

seq=26 sym=TSLA side=sell price=175.17 qty=1547
[FPGA->ARM] seq=26 action=BUY
  best_bid_px_1e4=1749900 best_bid_qty=2836
  best_ask_px_1e4=1750100 best_ask_qty=681
  spread_1e4=200 imbalance=2155

seq=27 sym=NVDA side=sell price=875.13 qty=3861
[FPGA->ARM] seq=27 action=BUY
  best_bid_px_1e4=8749800 best_bid_qty=3282
  best_ask_px_1e4=8750100 best_ask_qty=2259
  spread_1e4=300 imbalance=1023

seq=36 sym=AAPL side=sell price=185.75 qty=1795
[FPGA->ARM] seq=36 action=NOOP
  best_bid_px_1e4=1849900 best_bid_qty=939
  best_ask_px_1e4=1850300 best_ask_qty=818
  spread_1e4=400 imbalance=121
\end{lstlisting}
\fonte{Registro de execução dos autores na DE10-Nano}
\end{sourcecode}

A \autoref{tab:execucaoFuncionalDe10} relaciona cada conclusão funcional à evidência disponível. O ensaio não inclui ordem de saída, bolsa externa ou instrumento financeiro real.

\begin{table}[htb]
\caption{Evidências do ensaio funcional na DE10-Nano}
\label{tab:execucaoFuncionalDe10}
\begin{tabularx}{\textwidth}{lX}
\toprule
\textbf{Item} & \textbf{Evidência observada} \\ \midrule
Identificação & Assinatura, versão e geometria lidas antes dos eventos. \\
Comunicação & Conexão TCP local e sequências correlacionadas nas respostas. \\
Estado & Melhores preços, quantidades, \textit{spread} e desbalanceamento devolvidos. \\
Decisão & BUY e NOOP coerentes com os limiares nos exemplos registrados. \\
Limite & Não foram exercitados envio de ordens, risco ou conexão externa. \\ \bottomrule
\end{tabularx}
\fonte{Elaborado pelos autores a partir do registro do \autoref{codigo:saidaFpgaArm}}
\end{table}

\section{Resultados dos testes automatizados}\label{sec:validacaoTestes}

Os procedimentos do \autoref{quad:matrizValidacao} foram executados na versão revisada. Os dois testes C++ concluíram sem falha: o teste do invólucro e o teste rápido do núcleo de referência. Também concluíram os seis alvos VHDL: ponte básica, ponte com rajadas e retorno de índices, livro, estratégia gerada, motor e periférico Avalon-MM.

Os bancos de teste que instanciam a estratégia gerada emitiram avisos de \codigo{NUMERIC_STD} no instante inicial, antes de os sinais assumirem valores conhecidos, mas alcançaram suas asserções finais de sucesso. Isso indica uma oportunidade de inicialização mais limpa nos testes; não houve asserção funcional de falha. A evidência permanece limitada aos casos neles codificados.

O \autoref{quad:criteriosAceitacao} resume a quantidade de testes concluídos em cada grupo. Ele registra execução sem falha, e não porcentagem de cobertura do código.

\begin{tabframed}[htb]
\caption{Resultado da execução dos testes automatizados}
\label{quad:criteriosAceitacao}
\begin{tabularx}{\textwidth}{lXl}
\toprule
\textbf{Grupo} & \textbf{Escopo} & \textbf{Resultado} \\ \midrule
C++ & Invólucro MMIO e núcleo C++ de referência. & 2 de 2 testes concluídos. \\
Ponte VHDL & Registradores, filas, rajadas e contrapressão. & 2 bancos de teste concluídos. \\
Processamento VHDL & Livro, estratégia gerada e motor integrado. & 3 bancos de teste concluídos. \\
Avalon-MM & Acesso ao periférico pela interface de barramento. & 1 banco de teste concluído. \\ \bottomrule
\end{tabularx}
\fonte{Execução dos alvos \codigo{make cpp-test} e \codigo{make vhdl-test-all}}
\end{tabframed}

\section{Campanha quantitativa}\label{sec:benchmarkDesempenho}

Os valores desta seção foram transcritos de uma campanha com 1.000.000 de eventos, precedidos de 10.000 eventos de aquecimento, na DE10-Nano a 50~MHz. Como a saída JSON bruta não está versionada, a análise não pode recalcular estatísticas nem avaliar variação entre execuções independentes. Os números são reportados com esse limite.

A \autoref{tab:resultadosBenchmark} reúne as métricas registradas. A latência interna é medida pelo contador no FPGA. O tempo C++ mede somente o algoritmo de referência. A latência fim a fim mede o modo síncrono entre chamadas de envio e recepção no ARM. A vazão usa o modo em fluxo, com mais de um evento em trânsito.

\begin{table}[htb]
\caption{Valores registrados na campanha quantitativa}
\label{tab:resultadosBenchmark}
\begin{tabularx}{\textwidth}{lXl}
\toprule
\textbf{Métrica} & \textbf{Escopo} & \textbf{Valor} \\ \midrule
Latência interna FPGA & Comando aceito até resposta produzida. & 3 ciclos (60~ns) \\
Variação interna & Máximo menos mínimo do contador de ciclos. & Não resolvida; desprezível na resolução de 1 ciclo (20~ns) \\
Tempo médio C++ & Livro e regra no ARM, sem comunicação. & 315~ns/evento \\
Razão C++/FPGA & Razão reportada entre as médias não arredondadas. & 5,253 vezes \\
RTT médio & ARM--MMIO--FPGA--MMIO--ARM, modo síncrono. & 10.228~ns \\
RTT p99 & Percentil 99 do mesmo modo. & 10.310~ns \\
RTT máximo & Maior amostra registrada. & 1.146.360~ns \\
Amplitude do RTT & Máximo menos mínimo registrado. & 1.136.220~ns \\
Vazão & Modo em fluxo pelas filas MMIO. & 102.381 mensagens/s \\ \bottomrule
\end{tabularx}
\fonte{Campanha dos autores com \codigo{fpga_benchmark --mode full --messages 1000000 --warmup 10000}}
\end{table}

O programa calcula também mediana, mas esse valor não foi preservado no texto nem em arquivo bruto. Percentis 90 e 95 tampouco foram registrados. Acrescentá-los agora exigiria repetir a campanha na placa; não é possível derivá-los dos valores disponíveis.

\subsection{Variação interna e resolução do contador}\label{subsec:jitter}

O contador atribuiu três ciclos tanto ao mínimo quanto ao máximo para os eventos da campanha. Como sua resolução é de um ciclo, ou 20~ns, o experimento não distingue variações menores que esse intervalo. A variação interna é, portanto, tratada como desprezível na resolução adotada. Essa classificação vale apenas para o núcleo e para a configuração testados; ela não determina a variação de outra implementação, frequência, condição de espera ou estratégia regenerada.

O núcleo possui caminhos sincronizados e tamanhos fixos, o que explica a contagem constante nos casos medidos. Entretanto, o atendimento da frequência depende do relatório de temporização da síntese, e o sistema completo inclui componentes assíncronos em relação ao núcleo. Por isso, o texto não transfere a observação de três ciclos para o ARM, o MMIO ou a rede.

A comparação de 60~ns com 315~ns produz o fator de 5,253 para os núcleos isolados. Ela compara duas implementações do algoritmo sob o procedimento descrito, não FPGA e software como plataformas completas. Quando a comunicação é incluída, o RTT médio de 10.228~ns é cerca de 32,4 vezes o tempo médio do núcleo C++; portanto, o protótipo não reduz a latência da decisão quando usado pelo ARM dessa maneira.

\subsection{Comunicação e vazão}\label{subsec:gargaloMMIO}

Cada envio bem-sucedido exige, no mínimo, leitura dos dois ponteiros TX, escrita das oito palavras e publicação de \codigo{TX_HEAD}. A recepção exige leitura dos ponteiros RX, leitura das oito palavras e publicação de \codigo{RX_TAIL}. Consultas adicionais ocorrem enquanto não há espaço ou resposta. Assim, a diferença entre 60~ns internos e 10.228~ns de ida e volta é compatível com o custo acumulado de MMIO, espera ativa e execução no ARM, mas a campanha não decompõe esse custo por componente.

O máximo de 1.146.360~ns mostra que o caminho completo teve uma cauda muito maior que a média. Sem rastreamento do sistema operacional e do barramento, não é possível atribuir essa amostra exclusivamente ao escalonador Linux ou exclusivamente ao MMIO. A conclusão sustentada é mais restrita: a variabilidade aparece fora do contador do núcleo lógico.

A vazão de 102.381 mensagens/s corresponde a aproximadamente 6,55~MB/s de carga útil somando 32 bytes de evento e 32 bytes de resposta. Esse cálculo não representa todo o tráfego do barramento, pois exclui ponteiros e consultas repetidas. O valor superou o limiar de engenharia de 100.000 mensagens/s adotado pelos autores, mas o limiar não veio de uma exigência de bolsa e não permite afirmar adequação a HFT operacional.

\subsection{Resultado do experimento de cálculo energético}\label{sec:energiaBenchmark}

A \autoref{tab:energiaComparacao} aplica a \autoref{eq:energiaDecisao} às potências e aos tempos declarados na metodologia. O resultado é estimado, não medido. Em particular, a linha x86 combina duas hipóteses e não corresponde a uma execução realizada neste trabalho.

\begin{table}[htb]
\caption{Estimativa de energia por decisão sob as hipóteses adotadas}
\label{tab:energiaComparacao}
\begin{tabularx}{\textwidth}{lXXX}
\toprule
\textbf{Substrato} & \textbf{Potência adotada} & \textbf{Tempo adotado} & \textbf{Energia estimada} \\ \midrule
Lógica FPGA Cyclone~V & 50~mW & 60~ns, registrado & 3~nJ \\
Núcleo ARM Cortex-A9 & 1,5~W & 315~ns, registrado & 0,47~$\mu$J \\
Núcleo x86 ilustrativo & 15~W & 30~ns, hipótese & 0,45~$\mu$J \\ \bottomrule
\end{tabularx}
\fonte{Experimento de cálculo dos autores; potências e tempo x86 adotados como hipóteses}
\end{table}

Sob essas hipóteses, o valor calculado para a lógica FPGA é inferior aos valores calculados para os núcleos de processador. A diferença não deve ser interpretada como medição de eficiência da plataforma: os limites de potência não são equivalentes e o cálculo exclui ARM, memória, ponte MMIO, alimentação da placa e rede. O resultado indica apenas uma hipótese a ser verificada por instrumentação direta.

\section{Discussão frente aos trabalhos relacionados}\label{subsec:comparacaoLiteraturaBench}

A \autoref{tab:comparacaoNumerica} contextualiza os números sem tratá-los como ranking. As referências de FPGA incluem rede ou medem subsistemas com fronteiras diferentes; o protótipo mede ida e volta por uma ponte MMIO controlada pelo ARM. Frequência, dispositivo, protocolo e instrumentação também diferem.

\begin{table}[htb]
\caption{Contextualização das latências reportadas}
\label{tab:comparacaoNumerica}
\begin{tabularx}{\textwidth}{lXX}
\toprule
\textbf{Referência} & \textbf{Valor reportado} & \textbf{Fronteira principal} \\ \midrule
Este trabalho & 60~ns internos; 10.228~ns de RTT médio & Núcleo Cyclone V e ida e volta ARM--MMIO--FPGA. \\
\citeonline{Leber2011} & inferior a 1~$\mu$s & Processamento FPGA integrado ao caminho de rede descrito pelos autores. \\
\citeonline{Boutros2017} & 269~ns no subsistema; 869~ns de ida e volta & Kintex UltraScale a 156~MHz com HLS e interface de rede. \\
\citeonline{Lockwood2012} & inferior a 1~$\mu$s & Biblioteca FPGA com funções de rede e negociação. \\ \bottomrule
\end{tabularx}
\fonte{Elaborado pelos autores com dados de \citeonline{Leber2011}, \citeonline{Boutros2017} e \citeonline{Lockwood2012}}
\end{table}

Os valores da literatura não isolam o efeito de custo, ferramenta HLS ou dispositivo. A comparação permite apenas observar que retirar o processador de propósito geral do caminho crítico é uma decisão arquitetural recorrente nos trabalhos de menor latência. Implementar rede diretamente no FPGA seria outro projeto e permanece fora do escopo desta monografia.

\subsection{Limitações da evidência}\label{sec:limitacoesResultados}

Os testes demonstram os comportamentos codificados e o ensaio de placa demonstra integração no cenário registrado. A campanha fornece uma caracterização inicial, mas carece de repetições independentes, arquivo bruto versionado, medição de uso de CPU e decomposição do RTT. O fluxo de 5 mensagens/s não foi usado para determinar a vazão; o programa de medição usa eventos pré-construídos. Também não houve medição elétrica. A energia por decisão foi apenas calculada com potências de referência e limites físicos diferentes; por isso, não se conclui eficiência energética da plataforma nem se extrapolam os valores para uma instalação de mercado.
